\begin{table*}[hbtp]
  \caption{Optimal hyperparameters identified and the mAP value for each YOLO model after comprehensive Optuna tuning. The $mAP$ values are presented as \(\mu \pm \sigma\), where $\mu$ represents the mean and $\sigma$ represents the standard deviation across the 5-fold validation sets of D-Fire image dataset.}
  \label{tb:tune_results}
  \csvreader[
    separator=semicolon,
    no head,
    late after line=\\,
    late after last line=\\\hline,
    tabularray={
      width=\textwidth,
      colspec={|X[1.5,c]|X[c]|X[c]|X[c]|X[c]|X[c]|X[3.0,c]|},
      colsep = {0.15pt},
      hlines,
      vlines,
      cell{1}{1} = {r = 1}{valign = m}, % Network
      cell{1}{2} = {r = 1}{valign = m}, % Epochs
      cell{1}{3} = {r = 1}{valign = m}, % Learning rate
      cell{1}{4} = {r = 1}{valign = m}, % Batch size
      cell{1}{5} = {r = 1}{valign = m}, % Momentum
      cell{1}{6} = {r = 1}{valign = m}, % Weight decay
      cells = {font = \footnotesize, halign = c, valign = m},
      row{1} = {font = \bfseries\footnotesize},
      row{1} = {blue!50!gray!25, font = \bfseries\footnotesize},
      row{odd} = {blue!85!gray!7},
    }
  ]{4.table/tb_tune_results.csv}{}{
    \csvexpval\csvcoli & % Network
    \csvexpval\csvcolii & % Learning rate
    \csvexpval\csvcoliii & %Batch size
    \csvexpval\csvcoliv & % Momentum
    \csvexpval\csvcolv & % Weight decay
    \csvexpval\csvcolvi % avg. mAP
  }
\end{table*}

% Footnote text for the table
% \footnotetext[1]{The $mAP$ values are presented as \(\mu \pm \sigma\), where $\mu$ represents the mean and $\sigma$ represents the standard deviation across the 5-fold validation sets.}

